\documentclass[10pt,a4paper]{amsart}
\usepackage[margin=19mm]{geometry}
\usepackage{amsmath,amssymb,mathtools}
\usepackage[scr=rsfs,cal=euler]{mathalfa}
\usepackage{xfrac}
\usepackage[unicode,hidelinks,bookmarks=false]{hyperref}
\newcommand{\ie}{i.e.\ }
\newcommand{\cf}{cf.\ }
\newcommand{\resp}{respectively\ }
\newcommand{\Subsection}{\subsection}
\providecommand{\nobreakdash}{\nobreak\hbox{-}}
\setlength{\emergencystretch}{2em}
\multlinegap0pt
\usepackage{bm}
\usepackage{bbm}
\usepackage{dsfont}
\usepackage{xspace}
\usepackage{cancel}
\usepackage{mathtools}
\usepackage{amsthm}
\makeatletter
\@ifundefined{theorem}{\newtheorem{theorem}{Theorem}}{}
\@ifundefined{lemma}{\newtheorem{lemma}[theorem]{Lemma}}{}
\@ifundefined{remark}{\newtheorem{remark}[theorem]{Remark}}{}
\makeatother
\title[Uncertainty principles associated with the linear canonical ]{Uncertainty principles associated with the linear canonical Dunkl transform}
\author{Umamaheswari S, Sandeep Kumar Verma}
\date{Source: arXiv:2507.00530v1 (CC BY 4.0)}
\begin{document}
\maketitle

\noindent\emph{Source and licence.} Uncertainty principles associated with the linear canonical Dunkl transform. Version arXiv:2507.00530v1 (CC BY 4.0), distributed under the Creative Commons Attribution 4.0 International licence, as recorded in the arXiv licence metadata for this exact version and shown on the abstract page https://arxiv.org/abs/2507.00530v1. Full rights notice: licenses/arxiv-2507.00530-CC-BY-4.0.txt.\par\smallskip

\noindent\textit{Editorial scope.} Counted unit: the Theorem whose source label is \texttt{Th:3.4} in the article (source0.tex lines 357--364, source label \texttt{Th:3.4}), delivered together with the article's own complete original proof of it (source0.tex lines 365--406, one \texttt{proof} environment of 42 lines). No mathematics is written, repaired or re-derived: statement and proof are the article's own text line for line. Required context retained uncounted: le:3.1 (lines 233--237); 5.1 (lines 284--286). Every definition, display and label of those lines is retained unchanged. Omitted from this package and not redistributed: 1--232, 238--283, 287--356, 407--1058. Nothing inside a retained excerpt is omitted or altered: every retained body is delivered exactly as the article's lines are written, so no omission receipt of any kind accompanies this package. Citations: the retained excerpts carry no citation command at all, so no bibliography entry of the article is transcribed and no reference is redistributed. Reference adaptation: none was needed and none was made; every reference inside the delivered unit resolves inside it, so no unresolved reference or citation is produced. Printed slips kept literal: no printed word, symbol, hypothesis or number of the counted statement or of its proof was altered. Layout and class: the article's own class is \texttt{amsproc} with packages including amsmath, amsthm, amscd, amsfonts, amssymb, graphicx, color, subcaption, hyperref; the wrapper is the lane's pinned \texttt{amsart} editorial wrapper at 10pt on A4, which restates the theorem-like environments the retained text needs and adds the article's own macro definitions that the retained text needs (none), with extra wrapper packages (bm, bbm, dsfont, xspace, cancel, mathtools). The delivered numbering is the unit's own. Assets: no retained span contains a figure environment, a graphics-inclusion command, a PDF-page inclusion, a table or a TikZ picture; no image file is copied into this package, the package declares no asset dependency and no image file is redistributed with it. Licence: version arXiv:2507.00530v1 of this article is distributed under the Creative Commons Attribution 4.0 International licence, as recorded in the arXiv licence metadata for this exact version and shown on the abstract page https://arxiv.org/abs/2507.00530v1. A full rights notice is delivered at licenses/arxiv-2507.00530-CC-BY-4.0.txt. Characters: every retained excerpt is pure ASCII, so the delivered engine, XeLaTeX with its default Latin Modern text font, reports no missing character.
\begin{lemma}[Riemann-Lebesgue lemma] \label{le:3.1} If $f\in L_k^1(\mathbb{R})$, then  $D_k^M(f)$ belongs to $\mathcal{C}_0(\mathbb{R})$ and satisfies the following inequality 
\begin{equation*} \label{e:2.4}
    \|D_k^M(f)\|_{L_k^\infty(\mathbb{R})} \le \frac{1}{|b|^{k+1}}\, \|f\|_{L_k^1(\mathbb{R})}.
\end{equation*}
\end{lemma}

\begin{equation}\label{5.1}
 |f \chi_{{\mathbb{R}}\setminus(-r,r)}(y)|\le |y|^\alpha\,r^{-\alpha}\,|f(y)|.   
\end{equation}

 \begin{theorem}\label{Th:3.4}
Let  $f$ belongs to $ L
 _k^1(\mathbb{R})$ and $L^p_k(\mathbb{R}), \,\, 1<p\le 2$. Then 
 \begin{equation*}
 \|D_k^M(f)\|_{L_k^{q}(\mathbb{R})} \le C(b,q,k) \, \|f\|_{L_k^1(\mathbb{R})}^{\frac{qs}{2k+2+qs}}\, \||x|^s\,D_k^M(f)\|_{L_k^{q}(\mathbb{R})}^{\frac{2k+2}{2k+2+qs}}, \quad s>0,
 \end{equation*}
 where $q$ is the conjugate exponent of $p$ and $$C(b,q,k) = \left \{ 1+\frac{C_{k,b}^q}{2^{k+1}\Gamma(k+2)}\right\}^{\frac{1}{q }}.$$
 \end{theorem}

 \begin{proof}
    For  $r>0$, we consider
 \begin{equation} \label{e:4.12}
 \|D_k^M(f)\|^q_{L_k^{q}(\mathbb{R})} =    \|\left(1-\chi_{(-r,r)}\right) \, D_k^M(f)\|^{q}_{L_k^{q}(\mathbb{R})}+ \|\chi_{(-r,r)}\, D_k^M(f)\|^{q}_{L_k^{q}(\mathbb{R})}.   
 \end{equation}
By using the fact \eqref{5.1}, we have
 \begin{equation}\label{e:6.2}
 \|\left(1-\chi_{(-r,r)}\right) \, D_k^M(f)\|^{q}_{L_k^{q}(\mathbb{R})}\le \frac{\||x|^s\, D_k^M(f)\|^{q}_{L_k^{q}(\mathbb{R})}}{r^{qs}}.   
 \end{equation}
 On manipulating the second part of \eqref{e:4.12} and invoking
 Riemann-Lebesgue Lemma \ref{le:3.1}, we get
\begin{eqnarray*}
    \|\chi_{(-r,r)}\, D_k^M(f)\|^{q}_{L_k^{q}(\mathbb{R})} &\le& \|\chi_{(-r,r)}\|^{q}_{L_k^{q}(\mathbb{R})}\, \|D_k^M(f)\|^q_{L^\infty_k(\mathbb{R})} \\
    &\le&C_{k,b}^{q}\,\gamma_k\left((-r,r)\right)\,  \|f\|^{q}_{L_k^1(\mathbb{R})},
\end{eqnarray*}
where
\begin{eqnarray} \label{5.9}
\nonumber\gamma_k\left((-r,r)\right) &= &\int_{\mathbb{R}}  \chi_{(-r,r)}(x)\, d\mu_k(x)\\ 
&=& \frac{r^{2k+2}}{2^{k+1}\Gamma(k+2)}.
\end{eqnarray}
Thus, 
\begin{equation}\label{e:6.3}
    \|\chi_{(-r,r)}\,D_k^M(f)\|_{L^{q}_k(\mathbb{R})}^{q} \le C_{k,b}^{q}\,  \frac{r^{2k+2}}{2^{k+1}\Gamma(k+2)}\,\|f\|^{q}_{L_k^1(\mathbb{R})}.
\end{equation}
By combining  \eqref{e:4.12}, \eqref{e:6.2} and \eqref{e:6.3}, we obtain
\begin{equation*}
 \|D_k^M(f)\|_{L_k^{q}(\mathbb{R})}^{q} \le \frac{\||x|^s\, D_k^M(f)\|^{q}_{L_k^{q}(\mathbb{R})}}{r^{qs}}  \,  +  C_{k,b}^{q}\,  \frac{r^{2k+2}}{2^{k+1}\Gamma(k+2)}\,\|f\|^{q}_{L_k^1(\mathbb{R})}. 
\end{equation*}
By choosing $$r= \left[\frac{\||x|^s\, D_k^M(f)\|^{q}_{L_k^{q}(\mathbb{R})} }{\|f\|^{q}_{L_k^1(\mathbb{R})}}\right]^{\frac{1}{qs+2k+2}},$$ we deduce the required result
\begin{equation*}
    \|D_k^M(f)\|_{L_k^{q}(\mathbb{R})}\le C(b,q,k)\, \|f\|_{L_k^1(\mathbb{R})}^{\frac{qs}{2k+2+qs}}\, \||x|^s\, D_k^M(f)\|_{L_k^{q}(\mathbb{R})}^{\frac{2k+2}{2k+2+qs}}.
\end{equation*}
This completes the proof of the theorem.
\begin{remark}
    For $q=2$, the inequality in Theorem \ref{Th:3.4} reduces to \begin{equation*}
    \|f\|_{L_k^{2}(\mathbb{R})}\le C(b,k)\, \|f\|_{L_k^1(\mathbb{R})}^{\frac{s}{k+s+1}}\, \||x|^s\, D_k^M(f)\|_{L_k^{2}(\mathbb{R})}^{\frac{k+1}{k+s+1}}.
\end{equation*} 
\end{remark}



\end{proof}

\end{document}
